\documentclass{my_eth_beamer}

\usepackage{my_math_eth}

\numberwithin{equation}{section}

% -----------------------------------------------------------------------------
% this document
\title{Introduction to Survey Sampling}
\newcommand{\TITLE}{Stratified simple random sampling}
\newcommand{\NR}{4}
\subtitle{Part \NR: \TITLE}
\author[Tobias Schoch]{Tobias Schoch}
\date{2024}

% -----------------------------------------------------------------------------
\begin{document}

\setstretch{1.1}

\section{\TITLE}\label{sec:strat}
\setcounter{section}{\NR}
\setcounter{Bsp}{12}
\setcounter{Def}{10}
\setcounter{EinSatz}{1}


% -----------------------------------------------------------------------------
\begin{frame}
    \frametitle{Overview}
    \noindent \hspace{0.2em}\textcolor{fhnwblue}{\bfseries \NR~\TITLE}
    \begin{enumerate}[\NR.1]
        \item Estimation from sample data
        \item Statistical inference
        \item Sample size allocation
        \item Efficiency gains from stratification
        \item When and how to stratify?
    \end{enumerate}
\end{frame}

% -----------------------------------------------------------------------------
\begin{frame}
    \frametitle{\insertsectionnumber~\insertsection}
    \begin{itemize}
    \setlength{\itemsep}{14pt}
        \item The population $U$ is \alert{partitioned} into $H$ mutually
            exclusive (and exhaustive), nonempty \alert{strata}
            (subpopulations) $U_h$ ($h = 1, \ldots, H$).
        \item Stratum $U_h$ contains $N_h$ elements.
        \item It follows that
            \begin{equation*}
                U = \bigcup_{h = 1}^H U_h \qquad \text{and}
                    \qquad N = \sum_{h = 1}^H N_h,
            \end{equation*}

            \vspace{1em}

            \noindent or, in short notation, $U = \bigcup_h U_h$ and $N =
            \sum_h N_h$.
        \item The partition \alert{assigns} each element $i \in U$ to exactly
            one stratum.
    \end{itemize}
\end{frame}

% -----------------------------------------------------------------------------
\begin{frame}
    \begin{table}[htb]
        \caption{Population of farms in a U.S. county}
        \label{tab:strat_farm}
        {\small
        \centering
        \begin{tabular}{llrrr}
            \toprule
            \addlinespace
            Stratum & Farm area & Number of & Mean corn & Standard\\
            & (grouped) & farms & acreage & deviation \\
            \addlinespace
            \midrule
            1 &   0 -- 40  & 394 &  5.4 &  8.3\\
            2 &  41 -- 80  & 461 & 16.3 & 13.3\\
            3 &  81 -- 120 & 391 & 24.3 & 15.1\\
            4 & 121 -- 160 & 334 & 34.5 & 19.8\\
            5 & 161 -- 200 & 169 & 42.1 & 24.5\\
            6 & 201 -- 240 & 113 & 50.1 & 26.0\\
            7 & 241 --     & 148 & 63.8 & 35.2\\
            \midrule
            \multicolumn{2}{l}{Mean or total} & 2\,010 & 26.3 &--\\
            \bottomrule
            \addlinespace
            \multicolumn{5}{l}{\footnotesize{Source: \citet[][p.
                111]{cochran_1977}; areas in acres (1 acre $\approx 0.4$
                hectares).}}
        \end{tabular}
    }
    \end{table}
\end{frame}

% -----------------------------------------------------------------------------
\begin{frame}
    \begin{itemize}
    \setlength{\itemsep}{8pt}
        \item Stratification \alert{primarily} serves \alert{two purposes}:
            \vspace{0.25em}
            \begin{itemize}
            \setlength{\itemsep}{4pt}
                \item \alert{Efficiency:} more precise estimates than under \SRS
                    through suitable stratification and sample size allocation
                \item \alert{Estimation for planned domains:} subpopulations are
                defined as \emph{planned} domains with predetermined sample sizes
            \end{itemize}
        \item In practice, the allocation may also account for stratum-specific
            characteristics such as variability or anticipated response rates
            (Chapter \insertsectionnumber.\ref{sec:strat_zweck}).
       \item The \alert{notation} is somewhat \alert{cumbersome}. We therefore
           introduce the main notation for totals at this point.
        \item Values of the $y$-variable in stratum $h$ are identified using
            the \alert{double subscript} $ih$ and written as $y_{ih}$.
    \end{itemize}
\end{frame}

% -----------------------------------------------------------------------------
\begin{frame}
    \begin{itemize}
        \item The \alert{population total} of the $y$-variable \alert{in
            stratum} $h$ is denoted by $\tyh$ ($h = 1, \ldots, H$). We write:
        \vspace{-0.5em}
        \begin{align*}
            \Nh &= \text{size of stratum $h$,} r\\
            N &= \sum_{h = 1}^H \Nh \; = \; \text{population size,}\\
            \tyh &= \sum_{i \in U_h} y_{ih} \; = \;
                \text{$y$-total in stratum $h$,}\\
            \ty &= \sum_{h = 1}^H \tyh \; = \;
                \text{$y$-total in the population.}
        \end{align*}
        \item The \alert{population total} in can equivalently be written as
            \vspace{-0.5em}
            \begin{equation*}
                \ty = \sum_{h=1}^H \Nh \, \myh,
            \end{equation*}
    \end{itemize}
\end{frame}

% -----------------------------------------------------------------------------
\begin{frame}
    \hspace{1.75em} where $\myh$ is the \alert{mean} in stratum $h$, given by
    \begin{equation*}
        \myh = \frac{1}{\Nh}\sum_{i \in U_h} y_{ih} \qquad
            (h = 1, \ldots, H).
    \end{equation*}

    \begin{itemize}
    \setlength{\itemsep}{10pt}
        \item The \alert{population mean}, denoted by $\myu$, is the weighted
            sum of the stratum means $\myh$,
            \begin{equation*}
                \myu = \sum_{h = 1}^H W_h \, \myh,
            \end{equation*}
            \noindent where
            \begin{equation*}
                W_h = \frac{\Nh}{N} \qquad (h = 1, \ldots, H)
            \end{equation*}

            \vspace{1em}

            \noindent is a \alert{stratum weight}.
    \end{itemize}
\end{frame}

% -----------------------------------------------------------------------------
\begin{frame}
    \begin{itemize}
    \setlength{\itemsep}{12pt}
        \item The \alert{variance} of the $y$-variable in stratum $h$ is
            given by
            \begin{equation*}
                \varyh = \frac{1}{\Nh - 1} \sum_{i \in U_h} (y_{ih} - \myh)^2
                \qquad (h = 1, \ldots, H).
            \end{equation*}
        \item The \alert{population variance} is
            \begin{equation*}
                \varyu = \frac{1}{N - 1}\sum_{h=1}^H \sum_{i \in U_h}
                    (y_{ih} - \myu)^2.
            \end{equation*}
    \end{itemize}
\end{frame}

% -----------------------------------------------------------------------------
\begin{frame}
    \begin{Exkurs}[One-way analysis of variance]\label{exkurs:strat_anova}
        \vspace{-0.5em}
        An alternative expression for the variance $\varyu$ is
        \begin{equation*}
            \varyu = \frac{\SST}{N - 1},
        \end{equation*}

        \noindent where \alert{SST} denotes the total sum of squares. It can be
        decomposed into two components: the sum of squares within strata
        (\alert{SSW}) and the sum of squares between strata (\alert{SSB});
        formally,
        \begin{equation*}
            \SST = \SSW + \SSB,
        \end{equation*}
        \noindent with
        \begin{equation*}
            \SSW = \sum_{h = 1}^H (\Nh - 1) \varyh \quad \text{and} \quad
                \SSB = \sum_{h = 1}^H \Nh \big(\myh -\myu \big)^2.
        \end{equation*}
        \vspace{0.25em}
    \end{Exkurs}
\end{frame}

\subsection{Estimation from sample data}

% -----------------------------------------------------------------------------
\begin{frame}
    \frametitle{\insertsectionnumber.\insertsubsectionnumber~\insertsubsection}
    \begin{itemize}
    \setlength{\itemsep}{12pt}
        \item A \alert{stratified} probability sample has the following
            \alert{characteristics}:
        \vspace{0.25em}
        \begin{itemize}
        \setlength{\itemsep}{8pt}
            \item Each stratum $U_h$ is treated as a \alert{separate
                subpopulation}.
            \item A \alert{separate}, \alert{independent sample} $s_h$ of size
                $n_h$ is selected from each stratum $U_h$ ($h = 1, \ldots, H$).
        \end{itemize}
        \item Independence implies that
        \vspace{0.25em}
        \begin{itemize}
        \setlength{\itemsep}{8pt}
            \item a \alert{different} sampling design can, in principle, be
                used in each stratum;
            \item the \alert{variance} of an estimator can be
                \alert{decomposed} (in short notation) as
                \begin{equation*}
                    \var = \sum_h \mathsf{var}_{p,h}
                \end{equation*}
        \end{itemize}
    \end{itemize}
\end{frame}

% -----------------------------------------------------------------------------
\begin{frame}
    \begin{itemize}
    \setlength{\itemsep}{8pt}
        \item For \alert{practical reasons}, it is advisable to use designs
            from a single family (e.g., \SRS), particularly to keep variance
            estimation manageable.
        \item We \alert{restrict attention} to the case in which an \alert{SRS}
            is selected independently within each stratum.
        \item The \alert{sample sizes} $n_h$ ($h = 1, \ldots, H$) can be
            allocated across strata in \alert{different ways} (e.g.,
            proportional to $N_h$ or according to another rule).
        \item The stratified design is denoted by $p_{\,\mSTSI}$.
        \item The \alert{overall sample} $s$ and its size $n$ are obtained as
            the \alert{union} of the stratum samples $s_h$, $h = 1, \ldots, H$,
            and the sum of their sample sizes $n_h$,
            \begin{equation*}
                s = \bigcup_{h = 1}^H s_h \qquad \text{and}
                    \qquad n = \sum_{h = 1}^H n_h.
            \end{equation*}
    \end{itemize}
\end{frame}

% -----------------------------------------------------------------------------
\begin{frame}
    \begin{itemize}
    \setlength{\itemsep}{12pt}
        \item Let $s$ be a sample under design $p_{\,\mSTSI}$. The estimators
            of the \alert{mean} and \alert{variance} in stratum $h$ are
            \begin{equation*}
                \myhs = \frac{1}{n_h}\sum_{i \in s_h} y_{ih} \qquad \text{and}
                    \qquad \varyhs = \frac{1}{n_h - 1}\sum_{i \in s_h} (y_{ih}
                    - \myhs)^2 \qquad (h=1, \ldots,H),
            \end{equation*}
            \begin{itemize}
            \setlength{\itemsep}{12pt}
                \item \alert{provided} that $n_h > 0$ (mean) and $n_h > 1$
                    (variance), respectively.
                \item Both estimators are unbiased for the corresponding
                    population characteristics.
            \end{itemize}
        \item An unbiased \alert{estimator} of the \alert{population mean} is
            \begin{equation*}
                \mystrat = \sum_{h = 1}^H W_h \, \myhs.
            \end{equation*}
    \end{itemize}
\end{frame}

\setcounter{equation}{\the\numexpr\value{equation}+1\relax}

% -----------------------------------------------------------------------------
\begin{frame}
    \begin{itemize}
    \setlength{\itemsep}{12pt}
        \item The theoretical \alert{variance} of $\mystrat$ is
            \begin{equation}\label{eq:strat_var_mean}
                \var\big(\mystrat\big) = \sum_{h = 1}^H W_h^2
                    \left(1 - \frac{n_h}{\Nh}\right) \frac{\varyh}{n_h}.
            \end{equation}
        \item An unbiased \alert{variance estimator} is
            \begin{equation*}
                \widehat{\var}\big(\mystrat\big) = \sum_{h = 1}^H W_h^2
                    \left(1 - \frac{n_h}{\Nh}\right) \frac{\varyhs}{n_h}.
            \end{equation*}
        \item An unbiased estimator of the \alert{population total} $t_y$ is
            \begin{equation*}
                \tystrat = \sum_{h = 1}^H \Nh \, \myhs \qquad \text{or}
                    \qquad \tystrat = N \, \mystrat.
            \end{equation*}
    \end{itemize}
\end{frame}

% -----------------------------------------------------------------------------
\begin{frame}
    \begin{itemize}
    \setlength{\itemsep}{14pt}
        \item The \alert{theoretical} variance of $\tystrat$ is
            \begin{equation*}
                \var\big(\,\tystrat\big) = N^2 \,\var\big(\mystrat\big).
            \end{equation*}
        \item An unbiased \alert{variance estimator} is
            \begin{equation*}
                N^2 \; \varest\big(\mystrat\big).
            \end{equation*}
    \end{itemize}
\end{frame}

\subsection{Statistical inference}

% -----------------------------------------------------------------------------
\begin{frame}
    \frametitle{\insertsectionnumber.\insertsubsectionnumber~\insertsubsection}
    \begin{itemize}
    \setlength{\itemsep}{10pt}
        \item The design $p_{\mSTSI}$ of the overall sample consists of $H$
            stratum-specific designs, \alert{all based on \SRS}.
        \item It is therefore natural to construct \alert{confidence intervals}
            \alert{analogously} to those under \SRS (Chapter 3).
        \item There are \alert{two asymptotic settings}:
            \vspace{0.25em}
            \begin{enumerate}
            \setlength{\itemsep}{6pt}
                \item The number of strata $H$ is fixed and $n_h,N_h
                    \rightarrow \infty$.
                \item The stratum sizes $n_h$ and $N_h$ are fixed and $H
                    \rightarrow \infty$.
            \end{enumerate}
        \item A \alert{central limit theorem} can be established in both
            settings; see \citet{krewski_rao_1981}.
    \end{itemize}
\end{frame}

% -----------------------------------------------------------------------------
\begin{frame}
    \begin{itemize}
    \setlength{\itemsep}{10pt}
        \item For sufficiently large sample sizes, an approximate
            $100(1-\alpha)\%$ \alert{confidence interval} (CI) for $\myu$ is
            \begin{equation*}
                \mystrat \; \pm \; z_{1-\alpha/2}\; \ses(\mystrat),
            \end{equation*}
            where $z_{1-\alpha/2}$ is the $(1-\alpha/2)$ quantile of the
            $N(0,1)$ distribution.
            \vspace{0.75em}
            \begin{itemize}
            \setlength{\itemsep}{10pt}
                \item As under \SRS, this confidence interval is based on a
                    \alert{normal approximation}.
                \item Its accuracy depends on the finite populations within the
                    strata and on the stratum sample sizes $n_h$, in particular
                    when distributions are skewed.
                \item Since sampling within each stratum is \SRS,
                    \alert{Cochran's rule} from Chapter 3 can be used as a
                    \alert{rough guide within individual strata}.
            \end{itemize}
    \end{itemize}
\end{frame}

\subsection{Sample size allocation}\label{sec:strat_allocation}

% -----------------------------------------------------------------------------
\begin{frame}
    \frametitle{\insertsectionnumber.\insertsubsectionnumber~\insertsubsection}
    \begin{itemize}
    \setlength{\itemsep}{10pt}
        \item Let $n$ denote the \alert{predetermined} overall sample size.
        \item When \alert{planning} the survey, we must decide how to
            \alert{allocate} $n$ across the strata such that
            \begin{equation*}
                \sum_{h=1}^H n_h = n
            \end{equation*}

            \vspace{-0.5em}

            \noindent holds.

        \item We consider two methods for \alert{allocating} the total sample
            size $n$ across strata:
            \vspace{0.25em}
            \begin{itemize}
            \setlength{\itemsep}{6pt}
                \item \alert{Proportional} allocation
                \item \alert{Neyman--Tschuprow} (optimal) allocation
            \end{itemize}
    \end{itemize}
\end{frame}

% -----------------------------------------------------------------------------
\begin{frame}
    \frametitle{\normalsize Proportional allocation}
    \begin{itemize}
    \setlength{\itemsep}{8pt}
        \item Proportional allocation can be used when \alert{no information
            beyond} the stratum population sizes $N_1, \ldots, N_H$ is
            available.
        \item Under this allocation, the sample sizes $n_h$ are chosen
            \alert{proportional} to $\Nh$ ($h = 1, \ldots, H$).
        \item The stratum sample sizes are set according to
            \begin{equation*}
                n_h = \left[n\, \frac{\Nh}{N} \right] =
                    \big[n \, W_h \big], \quad \text{with}
                        \quad W_h = \frac{\Nh}{N} \qquad (h = 1, \ldots, H),
            \end{equation*}
            \noindent where $[x]$ denotes a \alert{rounding function} that
            rounds $x$ to an integer.
        \item Because the $n_h$ must be integers, rounding may result in
            $\sum_{h=1}^H n_h \neq n$. The rounded values must therefore be
            adjusted if the equality is to hold.
    \end{itemize}
\end{frame}

% -----------------------------------------------------------------------------
\begin{frame}
    \begin{itemize}
    \setlength{\itemsep}{8pt}
    \item The general \alert{variance formula} for the population mean under a
        stratified design in (\ref{eq:strat_var_mean}) simplifies for the
        \alert{strategy} $\big(p^{\,\mprop}, \mys^{\,\mstrat}\big)$ to
            \vspace{0.25em}
            \begin{equation*}
                \var^{\mprop}\big(\mys^{\,\mstrat}\big) = (1 - f)\frac{1}{n}
                    \sum_{h = 1}^H W_h \varyh.
            \end{equation*}
        \item We use the short notation $p^{\,\mprop}$ instead of
            $p_{\,\mSTSI}^{\,\mprop}$.
    \end{itemize}
    \vspace{1em}
    \begin{Bsp}[Corn acreage]
        \label{ex:strat_county_prop}
        A stratified \SRS of size $n=100$ is to be selected from $U$ (see
        Table \ref{tab:strat_farm}) using proportional allocation. The
        resulting stratum sample sizes $n_h$ are shown below.
    \end{Bsp}
\end{frame}

% -----------------------------------------------------------------------------
\begin{frame}
    \begin{Bsp_forts}[Cont.]
        {\small
        \vspace{0.5em}
        \begin{center}
        \begin{tabular}{lrrrrrrr}
            \toprule
            \addlinespace
            Stratum $h$ & 1 & 2 & 3 & 4 & 5 & 6 & 7\\
            \addlinespace
            \midrule
            $W_h = N_h/N \quad$ & 0.196 & 0.229 & 0.195 & 0.166 & 0.084 &
                0.056 & 0.074\\[0.5em]
            $n_h = [n W_h]$ & 20 & 23 & 19 & 17 & 8 & 6 & 7\\
            \addlinespace
            \bottomrule
        \end{tabular}
        \end{center}
        }
        \vspace{0.25em}
    \end{Bsp_forts}
    \vspace{1em}
    \noindent\textcolor{fhnwblue}{\bfseries Optimal allocation according to
    Neyman--Tschuprow}
    \begin{itemize}
    \setlength{\itemsep}{8pt}
        \item The allocation proposed by \citet{neyman_1934} and Tschuprow
            \citep[see][]{dalenius_1992} assumes that, \alert{in addition} to
            the stratum population sizes $N_1, \ldots, N_H$, the
            \alert{stratum-specific population variances} of the variable of
            interest, $\sigma_{1}^2, \ldots, \sigma_{H}^2$, are known.
   \end{itemize}
\end{frame}

% -----------------------------------------------------------------------------
\begin{frame}
    \begin{itemize}
    \setlength{\itemsep}{12pt}
        \item Under this allocation, the stratum sample sizes are set as
            follows:
            \begin{equation}\label{eq:strat_neyman}
                n_h = \left[n \, \frac{W_h \, \sigma_{h}}{\sum_{j=1}^H W_j \,
                \sigma_{j}}\right], \quad \text{with} \quad W_h = \frac{\Nh}{N}
                \qquad (h = 1, \ldots, H),
            \end{equation}
            \noindent where $[\cdot]$ denotes the rounding function.
        \item The allocation is \alert{variance-optimal} because it minimizes
            the variance in (\ref{eq:strat_var_mean}) subject to the constraint
            \begin{equation*}
                \sum_{h = 1}^H n_h = n.
            \end{equation*}
        \item The sampling design is denoted by $p^{\mopt}$ for short
            (full notation: $p_{\,\mSTSI}^{\,\mopt}$).
    \end{itemize}
\end{frame}

% -----------------------------------------------------------------------------
\begin{frame}
    \begin{itemize}
    \setlength{\itemsep}{14pt}
        \item The \alert{variance} of the estimator $\mys^{\,\mstrat}$ under
            the strategy $\big(p^{\mopt}, \mys^{\,\mstrat}\big)$ is
            \begin{equation*}
                \var^{\mopt}\big(\mys^{\,\mstrat}\big) =
                    \frac{1}{n}\left(\sum_{h = 1}^H W_h \sigma_h\right)^2 -
                    \frac{1}{N}\sum_{h = 1}^H W_h \varyh.
            \end{equation*}

        \vspace{0.25em}

        \item The Neyman--Tschuprow allocation has the following
            \alert{properties}:
            \vspace{0.25em}
            \begin{itemize}
            \setlength{\itemsep}{8pt}
                \item Strata with \alert{large population sizes} $\Nh$ tend to
                    receive \alert{larger} sample sizes $n_h$.
                \item Strata with \alert{large population variances}
                    $\sigma_h^2$ tend to receive \alert{larger} sample sizes
                    $n_h$.
                \item If $\sigma_{1} = \sigma_{2} = \ldots = \sigma_{H}$, the
                    allocation reduces to proportional allocation.
            \end{itemize}
    \end{itemize}
\end{frame}

% -----------------------------------------------------------------------------
\begin{frame}
    \begin{itemize}
    \setlength{\itemsep}{12pt}
        \item In practice:
            \vspace{0.5em}
            \begin{itemize}
            \setlength{\itemsep}{12pt}
                \item The \alert{population standard deviations} $\sigma_h$
                    ($h=1, \ldots, H$) are usually unknown. Estimates
                    $\widehat{\sigma}_h$ may be substituted for $\sigma_h$ in
                    (\ref{eq:strat_neyman}).
                \item \alert{Estimates} may be obtained, for example, from
                    \alert{historical data} or \alert{related surveys};
                \item Surveys usually collect many variables, whereas
                    Neyman--Tschuprow allocation is defined for a single
                    variable. It may therefore be useful to compute allocations
                    for several variables of interest and combine them into a
                    \alert{``compromise allocation.''}
            \end{itemize}
        \item If unit sampling costs differ across strata, the
            \alert{cost-optimal allocation} also depends on the
            stratum-specific costs.
    \end{itemize}
\end{frame}

% -----------------------------------------------------------------------------
\begin{frame}
    \begin{Bsp}[Corn acreage, cont. of Ex. \ref{ex:strat_county_prop}]
        \label{ex:strat_county_opt}
        \vspace{-0.25em}
        A stratified \SRS of size $n=100$ is to be selected from $U$ (see
        Table \ref{tab:strat_farm}) using variance-optimal allocation. The
        resulting stratum sample sizes $n_h$ are shown below.
        \begin{center}
            {\small
            \begin{tabular}{lrrrr}
                \toprule
                \addlinespace
                Stratum $h$ &\;$W_h = N_h / N$ & $\quad W_h\sigma_{h}$ &
                $\quad$ Share& $\quad n_h$\\
                \addlinespace
                \midrule
                1 & 0.196 & 1.627 & 0.096 & 10\\
                2 & 0.229 & 3.046 & 0.179 & 18\\
                3 & 0.194 & 2.929 & 0.172 & 17\\
                4 & 0.166 & 3.287 & 0.193 & 19\\
                5 & 0.084 & 2.058 & 0.121 & 12\\
                6 & 0.056 & 1.456 & 0.086 & 9\\
                7 & 0.074 & 2.605 & 0.153 & 15\\
                \midrule
                \multicolumn{2}{l}{Total} & & 1.000 & 100\\
                \bottomrule
            \end{tabular}
        }
        \end{center}
    \end{Bsp}
\end{frame}

\subsection{Efficiency gains from stratification}\label{sec:strat_efficiency}

% -----------------------------------------------------------------------------
\begin{frame}
    \frametitle{\insertsectionnumber.\insertsubsectionnumber~\insertsubsection}
    \begin{itemize}
    \setlength{\itemsep}{12pt}
        \item Strategies based on stratified sampling designs can, under
            certain conditions, achieve \alert{efficiency gains} relative to
            \SRS.
        \item We speak of an efficiency gain when an estimator has a
            \alert{smaller variance than the reference estimator}.
        \item All sampling designs considered have the \alert{same sample size}
            $n$.
        \item The \alert{reference estimator} is $\mys_{\,\mSRS} := \mys$,
            corresponding to the strategy $(p_{\,\mSRS}, \, \mys_{\,\mSRS})$ with variance
            $\var^{\mSRS}$.
        \item As before, $\var^{\mprop}$ and $\var^{\mopt}$ denote the
            variances under $\mSTSI$ with proportional and optimal allocation,
            respectively.
    \end{itemize}
\end{frame}

% -----------------------------------------------------------------------------
\begin{frame}
    \frametitle{\normalsize Proportional allocation vs. SRS}
    \vspace{-0.5em}
    \begin{itemize}
    \setlength{\itemsep}{8pt}
        \item The \alert{efficiency gain} of the strategy
            $(p^{\,\mprop},\, \mys^{\,\mstrat})$ over
            $(p_{\mSRS},\,\mys_{\,\mSRS})$ can be
            expressed as the \alert{difference in variances} (in shorthand
            notation),
            \begin{equation*}
                \var^{\mSRS} - \var^{\mprop} = \frac{1 - f}{n(N - 1)}\left\{\SSB
                    - \sum_{h = 1}^H (1 - W_h)\varyh \right\}.
            \end{equation*}
        \begin{itemize}
        \setlength{\itemsep}{8pt}
            \item[\goto] The strategy $(p^{\,\mprop},\,
                \mys^{\,\mstrat})$ is \alert{more efficient} than
                $(p_{\,\mSRS}, \, \mys_{\,\mSRS})$ if
                \begin{equation*}
                    \sum_{h = 1}^H (1 - W_h) \varyh < \SSB.
                \end{equation*}
            \item[\goto] Thus, proportional stratification is more efficient
                than \SRS when \alert{between-} \alert{stratum variation} is
                sufficiently large relative to \alert{within-stratum
                variation}.
        \end{itemize}
    \end{itemize}
\end{frame}

% -----------------------------------------------------------------------------
\begin{frame}
    \begin{itemize}
    \setlength{\itemsep}{8pt}
        \item The \alert{SSB} is \alert{maximized} when the strata are formed
            so that the stratum-specific \alert{means} $\myh$:
        \vspace{0.25em}
        \begin{itemize}
        \setlength{\itemsep}{6pt}
            \item are as \alert{different} as possible across strata,
            \item while the strata are as \alert{homogeneous} as possible
                internally.
        \end{itemize}
    \end{itemize}
    \vspace{0.75em}
    \noindent\textcolor{fhnwblue}{\bfseries Optimal vs. proportional
        allocation}
    \begin{itemize}
    \setlength{\itemsep}{8pt}
        \item The \alert{efficiency gain} of the strategy $(p^{\mopt},\,
            \mys^{\,\mstrat})$ over $(p^{\,\mprop},\, \mys^{\,\mstrat})$
            can be expressed as the following difference in variances:
            \begin{equation*}
                \var^{\mprop} - \var^{\mopt} = \frac{1}{n} \sum_{h = 1}^H W_h
                \big(\sigma_h - \widebar{\sigma}\big)^2,
            \end{equation*}
        \noindent where $\widebar{\sigma} = \sum_{h=1}^H W_h \sigma_h$ is the
            \alert{weighted mean of the standard deviations} $\sigma_h$.
   \end{itemize}
\end{frame}

% -----------------------------------------------------------------------------
\begin{frame}
    \begin{Bsp}[Corn acreage, cont. of Ex. \ref{ex:strat_county_prop} and
        \ref{ex:strat_county_opt}] \label{ex:strat_county_eff}
        \vspace{-0.25em}
        For the population of farms in Table \ref{tab:strat_farm}, we compare
        two strategies for estimating the population mean corn acreage:
        stratified \SRS with proportional allocation and with variance-optimal
        allocation. The estimator variances (rounded to two decimal places)
        are $\var\big(\mys^{\,\mprop}\big)=3.26$ and
        $\var\big(\mys^{\,\mstrat}\big)=2.73$. Thus, the strategy
        $(p^{\mopt}, \mys^{\,\mstrat})$ has a theoretical variance that is
        about 16\% smaller.
    \end{Bsp}
\end{frame}

\subsection{When and how to stratify?}\label{sec:strat_zweck}

% -----------------------------------------------------------------------------
\begin{frame}
    \frametitle{\insertsectionnumber.\insertsubsectionnumber~\insertsubsection}
    Reasons for using stratified sampling include:
    \vspace{0.25em}
    \begin{itemize}
    \setlength{\itemsep}{8pt}
        \item \alert{Efficiency considerations:}
            \begin{itemize}
            \setlength{\itemsep}{4pt}
                \item With suitable stratification and sample size allocation,
                    a stratified sample can yield \alert{more precise}
                    estimates.
                \item Gains in precision may arise when strata differ in their
                    \alert{means} or \alert{variances} and the allocation
                    exploits these differences.
                \item In particular, stratification is effective when strata
                    are internally \alert{homogeneous} with respect to the
                    variable of interest.
            \end{itemize}
        \item \alert{Estimation for subpopulations:}
            \begin{itemize}
            \setlength{\itemsep}{4pt}
                \item Stratification allows subpopulations to be defined as
                    \alert{planned domains} and their sample sizes to be set
                    to achieve the required precision.
            \end{itemize}
    \end{itemize}
\end{frame}

% -----------------------------------------------------------------------------
\begin{frame}
    \begin{itemize}
    \setlength{\itemsep}{12pt}
        \item \alert{Administrative convenience:} A national survey may, for
            example, be easier to organize if each region (e.g., canton)
            independently collects its assigned part of the sample. In this
            case, regions may provide a suitable stratification.
        \item \alert{``Balanced'' samples (control of sample composition):}
            \begin{itemize}
            \setlength{\itemsep}{6pt}
                \item A stratified sample can prevent a potentially
                    ``unbalanced'' sample composition. For example, an \SRS of
                    $n=200$ university students might contain substantially
                    more women (e.g., 130) than men (e.g., 70), even if the
                    population is evenly split by gender.
                \item This is not a theoretical problem, but the results may
                    appear more acceptable to a general audience if the
                    subsamples are ``balanced.''
        \end{itemize}
    \end{itemize}
\end{frame}

% -----------------------------------------------------------------------------

\begin{frame}
    \begin{itemize}
    \setlength{\itemsep}{8pt}
        \item In practice, stratified sampling requires that
            \alert{stratum membership is known before sampling}.
            \vspace{0.25em}
            \begin{itemize}
            \setlength{\itemsep}{4pt}
                \item[\goto] Stratification may therefore be theoretically
                    desirable but infeasible if the information needed to
                    assign units to strata is \alert{not available on the
                    sampling frame}.
            \end{itemize}
        \item If \alert{response rates} are expected to differ across strata,
            the initial sample sizes $n_h$ may be adjusted accordingly. Strata
            with lower expected response rates may, for example, be sampled
            more heavily.
        \item If only the stratum population sizes $N_h$, $h=1,\ldots,H$, are
            known, but unit-level stratum membership is unavailable before
            sampling, \alert{post-stratification} may be used instead.
    \end{itemize}
\end{frame}

% -----------------------------------------------------------------------------
\setcounter{section}{\the\numexpr\value{section}-1\relax}
\begin{frame}
    \frametitle{Literature}
    {\footnotesize
        \setstretch{1.1}
        \bibliographystyle{fhnw_en}
        \bibliography{master}
    }
\end{frame}

% -----------------------------------------------------------------------------
\lastpage

\end{document}
